\section{Lezione 11} \hfill 18/03/2026
\subsection{Aceri}
Gli aceri hanno un legno bianco indifferenziato, un pò marezzato e quasi setoso; sono usati per mobili, intarsi o strumenti musicali.
\subsection{Tiglio}
Il legno del tiglio è molto buono, perché è:
\begin{itemize}
    \item bianco indifferenziato;
    \item molto stabile;
    \item facilmente lavorabile e verniciabile (es. per porte, finestre).
\end{itemize}
Queste specie sono: 
\begin{itemize}
    \item mesofile (tollerano una certa quantità di ombra);
    \item vogliono suoli freschi/sciolti/profondi (piuttosto esigenti d'acqua);
    \item si trovano normalmente nel piano collinare o basso montano; 
    \item sono specie pioniere secondarie (si possono trovare nella faggeta in cui solitamente ci sono rari disturbi);
    \item naturalmente si trovano nelle forre (umidi e ricchi);
    \item la disseminazione è anemocora, con i semi alati;
    \item buona capacità pollonifera caulinare.
\end{itemize}
Queste specie ono state utilizzate in arboricoltura da legno (soprattutto in centroeuropa); inoltre, si trovavano spesso nei maggenghi, vicino alle cascine (in particolare il frassino). Con l'abbandono della mezza montagna, queste specie hanno invaso i pascoli e prati (in particolare da parte degli aceri e frassini).\\
Infatti, nelle Alpi ci sono molte aree di acero-frassineti di nuova formazione (alcune di 50-60 anni). Sono trovabili nell'area friulana vicino alla Slovenia, nel comasco, varesotto o nell'area a nord della Valbelluna. \\
Queste popolazioni sono effimere, ovvero prima o poi verranno sostituite da altre formazioni. Queste formazioni sono naturalmente fustaie; non possono essere ceduate (poichè sarebbe una conversione a ceduo, che è vietata in Italia).\\
Non è possibile lasciare tutti questi popolamenti come acero-frassineti, dato il limitato interesse di mercato. Tutte le zone senza interesse economico possono essere portate verso le definitive (es. faggio o rovere).\\
In queste generazione, quei popolamenti possono essere utilizzati, sia come selvicoltura di popolamento o d'albero.\\
Con la selvicoltura di popolamento si deve arrivare alle 300-500 piante ad ettaro: per fare ciò si devono eseguire con una serie di diradamenti.\\
Inizialmente, per fare le prove su questi nuovi popolamenti, sono stati fatti degli sfolli precoci, risultati poi un fallimento: ridurre subito il popolamento, formato da specie con una grande capacità di rinnovazione, porta il popolamento a ricostituirsi. In questi popolamenti è conveniente lasciare fare alla natura, e fare i diradamenti quando le piante sono alte 13-15 metri (il tempo di raggiungimento è diverso in base alla fertilità stazionale).\\
Il diradamento è sempre del tipo basso: lo si fa per rilasciare le piante migliori (è anche quello più semplice). Sarebbe però meglio fare un diradamento libero (o selettivo), cioè che va a scegliere i 300-500 candidati e li libera dall'eventuale competizione.\\
\fbox{ \centering
    \begin{minipage}{0.95\textwidth} 
I diradamenti possono essere fatti secondo diversi criteri: geometrico (non scelta) o di scelta.\\
I diradamenti geometrici solitamente vengono fatti sui popolamenti artificiali.\\
Il criterio di scelta viene utilizzato nei popolamenti naturali: può essere negativo o positivo (scelta delle piante brutte da eliminare o scelta delle piante belle da tenere). Il criterio negativo può essere dal basso, dall'alto o misto. Dal basso significa tagliare le piante perdenti dalla competizione, dall'alto significa eliminare piante brutte del piano dominante (es. biforcate). Il criterio positivo significa scegliere le piante belle da portare a fine turno e tagliare le piante che generano la competizione (che sono belle, fossero brutte non andrebbero in competizione).\\
Il criterio di bello delle latifoglie sono diversi rispetto a quelli delle conifere. Latifoglie: chioma molto larga, senza difetti, senza biforcazione.
    \end{minipage}
    }\\
    

Gli aceri frassineti hanno 1-2 diradamenti, a distanza di 5-10 anni uno dall'altro; possibilmente prima che le piante siano alte 20 metri. Dopo il diradamento si aspetta il fine turno per il taglio di rinnovazione.\\
Il diametro ... è intorno ai 30-40 cm. Questo obiettivo viene raggiunto in 40-50 anni. Il mercato straniero ha diametri obiettivi maggiori, 50-60 cm, con turno di 70-80 anni. Il problema è che a 80 anni l'acero non ha problemi, il frassino ha la caratteristica di tendere a creare un cuore nero (falso durame molto scuro, che si forma dentro la pianta): perciò va tagliato prima dei 50-55 anni\\
Il taglio di rinnovazione viene svolto con tagli successivi, lasciando una copertura del 50 percento, e se abbiamo le specie definitive, occorre lasciare alla naturale evoluzione. In alternativa occorre avere una nuova generazione di acero-frasineto.\\
Gli acero tiglieti sono più rari, il tiglio utilizzato è quasi tutto da importazione. Il tiglio, in fase giovanile, ha una corteccia simile a quella del .... i tiglieti ad alto fusto sono rari. \\
Quando ci sono tigli in bosco, è necessario allevare le piante piccole, in modo da fare del bel legname.\\
\subsection{Faggio}
Nel piendo del piano montano, regione esalpica ed appennina. La regione esalpica ha:
\begin{itemize}
    \item elevate precipitazioni;
    \item elevata umidità dell'aria; 
    \item temperature non troppo elevate, media intorno a 8-10 gradi;
    \item substrati sia calcarei che acidi che silicatici. 
\end{itemize}
Specie leader nella regione esalpica. Quando è nel suo optimum forma popolamenti puri. Ci sono pochi disturbi, poche perturbazioni, pochi sbalzi termici, poca neve dell'alta montagna. \\
Specie tipicamente oceanica, richiede una certa umidità dell'aria.\\
Sciafilo (come tutte le specie definitive), con un seme pesante (come tutte le definitive), apparato radicale... indifferente al substrato (sia suoli silicatici che basici, preferisce però i calcarei perchè più fertili), grande approfittatotre delle nebbie 9umidità atmosferica, sfruttabile grazie alla propria chioma che fa condensare e portabile alle radici mediante stem-flow. Questo compensa le ridotte precipitazioni.\\
Capacità pollonifera caulinare, non eccezionale, tende a formare popolamenti puri. ha una chioma molto densa e spessa, sotto la chioma di una faggeta pura mancano le altre specie. \\
Il tasso è più sciafilo del faggio; l'agrifoglio occupa la stessa nicchia ecologica del faggio in Sardegna.\\
Il faggio si può trovare naturalmente anche nell'area avanlpico-collinare e nella mesalpica. Nel collinare è possibile trovarlo nei colli euganei o nella collina di torino, nei versanti settentrionali: in queste aree non è troppo competitivo, ed è parte del popolamento misto di latifoglie.\\
Dove le condizioni sono migliori, ..., c'è la formazione di boschi misti. \\
Negli appennini, le faggete occupano la linea di confine degli alberi. \\
Il faggio è usato per legna da ardere. Viene considerata la migliore legna da ardere , 0.9-0.8 densitaa.... non scoppietta, gli impianti di combustione di biomassa vengono testati con il legno di faggio. Viene usato per gli stecchi del gelato magnum. ha delle lenticelle. utilizzato per i pannelli multistrato, oggettistica e soprattutto mobili che devono sopportare forti sollecitazioni (molto robusto). È il legno ideale per fare le doghe dei letti, tavoli e sedie.\\
L'areale del faggio è centro-europeo, dalla Sicilia, lungo tutto l'Appennino, tutte le Alpi, fino ai Pirenei, a Nord si arriva alla Svezia meridionale e fino ai carpazi. In Inghilterra ci sono molti faggi, ma è una specie indigena.\\
Non viene usato per tutto ciò che è a contatto per l'esterno , non viene usato per le travature. Il faggio, a contatto con l'acqua, marcisce molto velocemente. \\
È possibile allevarlo sia a ceduo che a fustaia.\\
Allevamento a ceduo. Ha una discreta capacità pollonifera, al primo anno i polloni non crescono bene, è una specie ad accrescimento lento, perde la capacità pollonifera intorno ai 40-50 anni. In base alle condizioni ecologiche, perde la capacità pollonifera prima se le condizioni sono buone.\\
I botanici riconoscono 3 ..... di faggio: luzulo fagion (luzula nivea, graminea che fa una pannocchietta bianca, su substrati acidi), cephalanteron fagion (cephalantera, orchidea su suoli basici) e eufagion.\\
Nell'eufagion il faggio perde la capacità pollonifera prima, perché più fertile.\\
Il ceduo è sempre stato un quasi obbligo, data la fame di energia sempre presente in Italia. La ceduazione avveniva nelle zone pìù ricche, con un ceduo matricinato, con 60-80 matricine ad ettaro e turni compresi tra i 20 e 30 anni (relativamente lunghi).\\
È una specie a seme pesante, che comincia a produrlo in abbondanza dopo i 30 anni: la ceduazione non permette il rilascio del seme. Le matricine erano 2/3 dell'età del seme ed 1/3 erano dell'età del ceduo. Molto spesso, venivano rilasciate sempre le stesse matricine (piante lupo): oggi queste piante sono interessanti per la biodiversità.\\
Dal faggio si produceva anche il carbone vegetale: a parità di potere calorifico, pesa 6 volte in meno rispetto alla legna. I cedui di faggio erano disseminati di carbonaie.  ginestra dei carbonai. Il carbone vegetale, oltre per la combustione, viene usato per i filtri e in farmacopea. Un vecchio ceduo è individuabile dall'area carbonile.\\
La pasciona del faggio è ogni 5 anni e c'è una semipascione di 3-4 anni.\\
Alcune zone, tipo in Veneto e Friuli, c'era una maggiore richiesta di energia: è stato creato il ceduo a sterzo, ovvero un ceduo disetaneo. Sulla stessa ceppaia, c'erano polloni giovani, medi e vecchi. Se il turno era di 30, si entrava in bosco 3 volte, ogni 10 anni: in questo modo c'erano 3 generazioni di polloni per ogni ceppaia. L'aspetto positivo è che l'entrata in bosco poteva avvenire più frequentemente. D'altro canto, il spollonamento non doveva danneggiare i polloni presenti in ceppaia; inoltre i polloni rimasti entravano in competizione con quelli nuovi che devono uscire (la produzione era minore rispetto a ceduo matricinato). Il continuo sterzo comporta un innalzamento delle ceppaie (fino ad un metro). Oggi lo sterzo non è più svolto.\\
Sarebbe sempre meglio avere una fustaia, perché possiede tutti gli assortimenti che si vogliono. La fustaia richiede più anni per crescere e maturare (come viene svolta in Francia e Germania).\\
Non è consigliato svolgere il taglio raso, perché é sciafilo: cresce meglio sotto l'ombra (ovvero faggio su faggio). Un taglio raso in faggeta farebbe crescere prima tutte le altre specie.\\
I popolamenti di faggio sono tendenzilamente coetaneiformi. Neanche il taglio salturario si presta bene a questo popolamento.....Conviene fare i tagli successivi: forma di trattamento per la creazione di una fustaia monoplana/coetaneiforme. Sono costituiti da 2 tagli principali e due tagli che possono esistere o non esistere. I primi due sono tagli si sementazione e sgombero. Gli altri due tagli sono secondari (tra sementazione e sgombero) e quello di preparazione (una decina d'anni prima della sementazione). 